\section{Uncertainty quantification: constructing the covariance blocks}
\label{sec:systematics}



The uncertainty budget combines three independent blocks --- systematic
(universe), statistical (bootstrap), and ML (training-seed variation) --- summed assuming
independence (distinct RNGs / physics sources). The historical 3D budget (\S\ref{sec:3d-syst}) was built as its analog on the
3D grid; no 4D uncertainty magnitude is reported (\S\ref{sec:4d}). Full mathematical procedures, with worked examples
on the reported numbers, are in App.~\ref{app:statmethods}. This section describes
how each block is built; the resulting budget, and its agreement with the published
total, are reported with the measurement in \S\ref{sec:results-uq}.


\subsection{Systematic universes}
\label{sec:uq-universes}
The detector and model systematics use the MAT-MINERvA universe machinery: 44
error bands and 187 universes in total, identical to the published analysis's
\texttt{GetStandardSystematics} set. Each universe is re-unfolded with the
central estimator seed. All per-band covariances use the MAT convention:
universe-mean centering and biased $1/N$. For a $\pm1\sigma$ pair this is the
outer product of the half-difference; a common displacement of both endpoints
from CV is excluded from the variance. The 100 PPFX
universes use the same mean-centered $1/N$ formula.

Top 2D bands: Flux \SI{\fluxBand}{\percent}, Muon\_Energy\_MINOS
\SI{2.31}{\percent}, MinosEfficiency \SI{1.47}{\percent},
Muon\_Energy\_MINERvA \SI{1.29}{\percent}. Figure~\ref{fig:uqbands} shows the
grouped fractional-uncertainty projections in the published Fig.~6/7 style
(the full $205\times205$ covariance projected exactly, $C_{1D}=PC_{2D}P^{\top}$,
before dividing by the 1D central value).

\begin{figure}[H]
  \centering
  \includegraphics[width=0.49\textwidth]{MEFHC_fig6_7_uncertainty_pt}
  \includegraphics[width=0.49\textwidth]{MEFHC_fig6_7_uncertainty_pz}
  \caption{Fractional uncertainty on the 2D cross section projected onto $\pT$
  (left) and $\pPar$ (right), grouped by source in the style of Figs.~6--7 of
  Ref.~\cite{MINERvA:2021owq}: flux, interaction models, normalization,
  hadronic response, muon reconstruction, and the statistical band.}
  \label{fig:uqbands}
\end{figure}

\paragraph{Flux $1/\Phi$ normalization.}
The dominant flux term enters as $\sigma\propto1/\Phi$; varying $\Phi$ per flux
universe, rather than dividing every universe by the CV $\Phi$, gives a Flux
band of \SI{\fluxBand}{\percent} (versus $\sim\SI{1}{\percent}$ under the CV
convention) and a standalone combined budget of \SI{\uqMedian}{\percent}. The \SI{1.4}{\percent} target-nucleon
normalization is a fully-correlated rank-1 band. The 100-universe flux
covariance has a participation-ratio effective rank of $\approx1$ (one coherent
normalization mode carries \SI{99.6}{\percent} of the variance), and its leading
mode correlates at $\rho=\rhoFluxMuon$ with the rank-1 muon-energy-scale band:
the rederived Bashyal~\cite{Bashyal:2021tqu} joint block, whose cross term is
rank-2, with two independent estimators giving $\rho=0.87$--$0.90$. Adding the
block leaves the systematic rank unchanged --- its directions already lie within
the span of the universe covariance (App.~\ref{sec:rank}).

\subsection{Statistical uncertainty: Poisson bootstrap}
\label{sec:uq-boot}
The statistical block is a 300-replica per-event Poisson(1) weight bootstrap
applied jointly to data and simulation, with two independent sub-RNGs (data
and MC are independent variance sources), the MC reco and truth weights riding
a \emph{single} per-event draw (preserving the within-event migration
correlation), and the GBDT seed pinned per replica so the bootstrap and
training-seed variances are separable. Each replica is a full re-unfold; the
covariance is the sample covariance over replicas
(Fig.~\ref{fig:bootstrap}). Our bootstrap block is genuinely smaller than the
published statistical block ($\approx2.55\times$ in square-root trace). We report this as an
\emph{observed} difference in the statistical covariance rather than a
demonstrated causal efficiency advantage; the data-only stream of the data/MC split
bootstrap ($C_{\mathrm{data}}/\texttt{StatOnly}=0.356$) shows it is not a
missing MC term, but cannot separate an efficiency gain from regularization bias
or undercoverage (App.~\ref{sec:rank}).

\begin{figure}[H]
  \centering
  \includegraphics[width=0.49\textwidth]{uq_spread_2d}
  \includegraphics[width=0.49\textwidth]{uq_corr_2d}
  \caption{300-replica Poisson bootstrap: per-bin relative spread of the
  unfolded 2D cross section (left) and the bootstrap correlation matrix on the
  205 reported bins (right).}
  \label{fig:bootstrap}
\end{figure}

\subsection{ML/optimization uncertainty: training-seed variation}
\label{sec:uq-ml}
ML stochasticity is budgeted by re-unfolding with 10 independent estimator
seeds at fixed data and fixed universe (Fig.~\ref{fig:seedscan}); the sample
covariance over seeds is the ML block. It is the smallest block
(median \SI{0.166}{\percent} per bin). Two refinements harden it: (i) a
\emph{train/test-split} variant, which also re-draws the internal training
split, gives a band $1.24\times$ the pure-seed band
--- the 2D budget of Table~\ref{tab:uq} uses the pure-seed band; (ii) the ensemble-mean central value (averaging the
seed ensemble) agrees with the production single-run CV to \SI{0.28}{\percent}.
\begin{figure}[H]
  \centering
  \includegraphics[width=0.49\textwidth]{seedscan_spread_2d}
  \includegraphics[width=0.49\textwidth]{seedscan_band_pt}
  \caption{ML training-seed variation: per-bin relative spread over 10 estimator seeds (left)
  and the resulting fractional band projected onto $\pT$ (right).}
  \label{fig:seedscan}
\end{figure}

\subsection{Combination and the unified-throw test}
\label{sec:uq-combine}
The nominal covariance decomposition is the block sum
$C = C_{\mathrm{syst}} + C_{\mathrm{stat}} + C_{\mathrm{ML}}$, which assumes
the blocks are independent \emph{and} that the unfolding responds linearly to
simultaneous perturbations (so that per-band covariances add). The second
assumption is tested directly by the \emph{unified throw}
(App.~\ref{sec:universes}; \S\ref{sec:eavailw-construction}): joint re-unfolds with all reweightable knob bands and the
flux multisim perturbed simultaneously per event.  The corrected construction
uses the actual $-1\sigma$ and $+1\sigma$ event-weight endpoints, interpolates
asymmetrically between them, holds the unfolding-estimator seed fixed for every
throw, and forms the joint covariance about the throw mean.  The joint
displacement of that mean from the CV is stored separately as a per-bin vector
and its norm; per-band endpoint displacements are not stored. The matched block comparator
uses the MAT convention: universe-mean centering and the biased $1/N$
normalization, with both endpoints re-unfolded. A single-throw superposition
probe is insufficient to test the nonlinear response. The construction does not
subtract a scalar run-to-run jitter floor, and the 5D result uses its strict
joint ensemble and matched-endpoint comparator.
% Implementation: nd-unfolding/unified_throw_cov.py

More generally, this reframes the block sum as an approximation --- valid only
where the unfolding responds linearly --- rather than as a default
construction.  We retain the per-source block decomposition for attribution,
but use it only after joint throws and matched block endpoints pass the stated
validation tests.

\subsection{What changes for the higher-dimensional products}
\label{sec:uq-higherdim}
Dump-all shifted branches are limited to CV event-loop support. For the five
genuinely kinematic bands (BeamAngleX/Y, MuonResolution, and
Muon\_Energy\_MINERvA/MINOS), selection-complete active-universe event loops
are the production migration gate. That replacement is in place for the adopted
five-axis covariance, whose lateral block is built from those selection-complete
bands; higher-dimensional products still built from the dump-all lateral
branches remain candidates. Replacing the support-limited lateral block by the
selection-complete one moves that block's square-root trace by
\SI{\latTraceMove}{\percent}. That is a change in a square-root trace and is
\emph{not} a bound on any individual bin: across the \latBins\ reported bins
the per-bin ratio of active to support standard deviation runs \latRatioMin\ to
\latRatioMax, and within the top \SI{1}{\percent} of bins by support variance
--- which carry \SI{\latTopVarShare}{\percent} of the total --- it runs
\latRatioTopMin\ to \latRatioTopMax. A small share of the aggregate variance
does not make a bin irrelevant to a quoted result.

\paragraph{Scope of the selection-complete replacement.} The selection-complete
replacement covers the five detector bands implemented through shifted
reconstructed kinematics. MINOS efficiency and the three GEANT
hadron-interaction bands (GEANT\_Neutron, GEANT\_Pion, GEANT\_Proton) retain
their weight-based treatment. This classification establishes their exclusion
from the kinematic replacement; it does not independently validate the
completeness of the hadronic-response model for the $\Eavail$ and $W$
measurement.

For the corrected 5D construction, the ML ensemble fixes the estimator seed
at 42 and varies the train/test split seed. It therefore measures split response,
not the full estimator stochasticity. A dedicated fixed-data estimator-seed scan
(twelve estimator seeds at fixed data and fixed split) closes this scope: its
covariance matrix has $\sqrt{\mathrm{Tr}\,C}=\SI{\gbdtAiEstTrace}{cm^2/nucleon}$,
\SI{\gbdtAiEstFrac}{\percent} of the candidate split-varied band
(\SI{\gbdtMlSplitTrace}{cm^2/nucleon}), so the adopted band is dominated by
genuine estimator stochasticity rather than by the particular train/test
partition. The scan remains an auxiliary robustness check rather than an
independent matrix --- split and estimator responses are not assumed
uncorrelated --- and both bands are minor components, about three percent of
the combined budget in $\sqrt{\mathrm{Tr}\,C}$.

For an earlier 5D candidate the unified-throw construction of
\S\ref{sec:uq-combine} was carried out both mean-centered
and CV-centered; no magnitude from either is quoted, and neither is adopted. The
adopted product is a later, separate \texttt{cv}-variant build
(\S\ref{sec:eavailw-construction}; App.~\ref{app:history-5dcandidates}).
% Routed record: App.~\ref{app:provenance} (\texttt{gbdtfive-macros}).

Repeating all 187 universe unfolds and the CV with per-universe background
subtraction changes the combined \emph{block-sum} sqrt-trace by only
\SI{0.28}{\percent} and the \emph{adopted} covariance by \SI{0.18}{\percent},
a numerically negligible bound on the background-CV choice (App.~\ref{app:history}).
% (\SI{0.18}{\percent} is after the flux (J28) correction; App. H keeps the
% superseded-product wording and the inflation-damping explanation.)
The 4D and extended-fiducial results do not yet have covariances
constructed with this procedure; the 5D covariance is not transferred to them.

% The chi^2 caveat formerly here moved to sec_results.tex (\S\ref{sec:results-compare}).
